\documentclass[11pt,twoside,openright,spanish]{report}
%\documentclass[runningheads,a4paper]{llncs}
%\documentclass{article}
\usepackage{array}
\newcolumntype{L}{>{\centering\arraybackslash}m{3cm}}
\usepackage{booktabs}
\usepackage[utf8]{inputenc}
\usepackage[spanish,es-tabla,mexico]{babel}
\usepackage{UNAMThesis}
\usepackage{titlesec, blindtext, color}
\usepackage{amsmath}
\usepackage{amsfonts}
\usepackage{fancyhdr}   
\usepackage{tensor}
\usepackage{multirow} 
\usepackage{tabu}
\usepackage[style]{fncychap}
\usepackage{tabularx}
\usepackage{makeidx}
\usepackage{emptypage}
\usepackage{floatrow}
%\usepackage{\tiny }{mathrsfs}
\usepackage[font=footnotesize,labelfont=bf]{caption} % Pie de figura
\usepackage{tabulary}
\usepackage[top=0.8in]{geometry} % Margen superior más adecuado
\usepackage{makeidx}
\usepackage{verbatim}
\usepackage{float} % Permite poner las imágenes en donde queramos
\usepackage{amssymb}
\usepackage{subfig}
\usepackage{layout}
\usepackage{calligra} 
\usepackage[dvipsnames,table,xcdraw]{xcolor}
\usepackage{graphicx} % Nos permite utilizar imágenes.
\usepackage{graphics}
\usepackage{caption}
\usepackage{mwe}
\usepackage{hyperref}
\usepackage{framed}
\usepackage{leftidx}
\usepackage{dsfont}
\usepackage{mathtools}
\usepackage{enumerate}
\usepackage{chngcntr} % Para no resetear notas de pie de página
\usepackage{setspace}
\usepackage{epstopdf}
\usepackage{setspace}
\usepackage{booktabs}
\usepackage{url}
\usepackage{calc}  
\usepackage{calrsfs}
\usepackage{enumitem}
\usepackage[T1]{fontenc}
\usepackage[bitstream-charter]{mathdesign}
\usepackage{etoolbox}
% code listing settings
\usepackage{listings}
\usepackage{wrapfig,lipsum,booktabs}
\usepackage[customcolors,shade]{hf-tikz}
\usepackage[authoryear,round]{natbib}


 

\DeclareFloatVCode{myrowsep}{\vskip 4ex}

\newcounter{bibcount}
\makeatletter
\patchcmd{\@lbibitem}{\item[}{\item[\hfil\stepcounter{bibcount}{\thebibcount.}}{}{}
\setlength{\bibhang}{2\parindent}
\renewcommand\NAT@bibsetup%
[1]{\setlength{\leftmargin}{\bibhang}\setlength{\itemindent}{-\parindent}%
	\setlength{\itemsep}{\bibsep}\setlength{\parsep}{\z@}}
\makeatother


\bibliographystyle{apalike2mod}

\hfsetbordercolor{blue}

\captionsetup[table]{position=bottom}

\apptocmd{\thebibliography}{\csname cleardoublepage phantomsection\endcsname\addcontentsline{toc} {chapter}{Bibliografía}}{}{}

\makeatletter
\patchcmd{\ttlh@hang}{\parindent\z@}{\parindent\z@\leavevmode}{}{}
\patchcmd{\ttlh@hang}{ }{}{}{}
\makeatother

\numberwithin{equation}{chapter}
\numberwithin{figure}{chapter}
\numberwithin{table}{chapter}
\counterwithout{footnote}{chapter}

\renewcommand{\arraystretch}{1.2}

%Interlineado
\renewcommand{\baselinestretch}{1.5}

\setlength{\headheight}{15pt} 

\logounam{Imagenes/Escudo-UNAM}
\logoinstitute{Imagenes/Escudo-IBT}
\pagenumbering{roman}
\flushbottom
\newtheorem{theorem}{Theorem}
\newtheorem{acknowledgement}[theorem]{Acknowledgement}
\newtheorem{algorithm}[theorem]{Algorithm}
\newtheorem{axiom}[theorem]{Axiom}
\newtheorem{case}[theorem]{Case}
\newtheorem{claim}[theorem]{Claim}
\newtheorem{conclusion}[theorem]{Conclusion}
\newtheorem{condition}[theorem]{Condition}
\newtheorem{conjecture}[theorem]{Conjecture}
\newtheorem{corollary}[theorem]{Corollary}
\newtheorem{criterion}[theorem]{Criterion}
\newtheorem{definition}[theorem]{Definition}
\newtheorem{example}[theorem]{Example}
\newtheorem{exercise}[theorem]{Exercise}
\newtheorem{lemma}[theorem]{Lemma}
\newtheorem{notat}[theorem]{Notat}
\newtheorem{problem}[theorem]{Problem}
\newtheorem{proposition}[theorem]{Proposition}
\newtheorem{remark}[theorem]{Remark}
\newtheorem{solution}[theorem]{Solution}
\newtheorem{summary}[theorem]{Summary}
\newenvironment{proof}[1][Proof]{\textbf{#1.} }{\ \rule{0.5em}{0.5em}}

\newcommand{\grad}{\hspace{-2mm}$\phantom{a}^{\circ}$}



\renewcommand{\sin}{\operatorname{\sen}}

\lfoot[]{}
\cfoot[]{}
\rfoot[]{}
\renewcommand{\footrulewidth}{0pt}
\renewcommand{\headrulewidth}{0.1pt}

%%%%%%%%%%%%%%%%%%%% Reemplazamos l con elle %%%%%%%%%%%%%%%%%%%%%%
\mathcode`l="8000
\begingroup
\makeatletter
\lccode`\~=`\l
\DeclareMathSymbol{\lsb@l}{\mathalpha}{letters}{`l}
\lowercase{\gdef~{\ifnum\the\mathgroup=\m@ne \ell \else \lsb@l \fi}}%
\endgroup

%%%%%%%%%%%%%%%%%%%%% Norma y valor absoluto %%%%%%%%%%%%%%%%%%%%%%
\DeclarePairedDelimiter\abs{\lvert}{\rvert}%
\DeclarePairedDelimiter\norm{\lVert}{\rVert}%

% Swap the definition of \abs* and \norm*, so that \abs
% and \norm resizes the size of the brackets, and the 
% starred version does not.
\makeatletter
\let\oldabs\abs
\def\abs{\@ifstar{\oldabs}{\oldabs*}}
\let\oldnorm\norm
\def\norm{\@ifstar{\oldnorm}{\oldnorm*}}
%%%%%%%%%%%%%%%%%%%%%%%%%%%%%%%%%%%%%%%%%%%%%%%%%%%%%%%%%%%%%%%%%%%
% Change Colors https://en.wikibooks.org/wiki/LaTeX/Colors
\hypersetup{
	bookmarks=true,         % show bookmarks bar?
	unicode=true,          % non-Latin characters in Acrobat’s bookmarks
	pdftoolbar=true,        % show Acrobat’s toolbar?
	pdfmenubar=true,        % show Acrobat’s menu?
	pdffitwindow=false,     % window fit to page when opened
	pdfstartview={FitH},    % fits the width of the page to the window
	pdftitle={Tesis Uzmar},    % title
	pdfauthor={Uzmar Gómez},     % author
	pdfsubject={Subject},   % subject of the document
	pdfcreator={Uzmar Gómez},   % creator of the document
	pdfproducer={Uzmar Gómez}, % producer of the document
	pdfkeywords={}, % list of keywords
	pdfnewwindow=true,      % links in new PDF window
	colorlinks=true,       % false: boxed links; true: colored links
	linkcolor=BlueViolet,          % color of internal links (change box color with linkbordercolor)
	citecolor=BrickRed,        % color of links to bibliography
	filecolor=NavyBlue,      % color of file links
	urlcolor=NavyBlue           % color of external links
}


\newenvironment{changemargin}[3]{
	\begin{list}{}{
			\setlength{\topsep}{#3}
			\setlength{\leftmargin}{#1}
			\setlength{\rightmargin}{#2}
			\setlength{\listparindent}{\parindent}
			\setlength{\itemindent}{\parindent}
			\setlength{\parsep}{\parskip}
		}
		\item[]}{\end{list}}	
	
\usepackage{listings}
\usepackage{color}

\definecolor{dkgreen}{rgb}{0,0.6,0}
\definecolor{gray}{rgb}{0.5,0.5,0.5}
\definecolor{mauve}{rgb}{0.58,0,0.82}

\lstset{frame=tb,
	language=R,
	aboveskip=3mm,
	belowskip=3mm,
	showstringspaces=false,
	columns=flexible,
	basicstyle={\small\ttfamily},
	numbers=left,
	numberstyle=\tiny\color{gray},
	keywordstyle=\color{blue},
	commentstyle=\color{dkgreen},
	stringstyle=\color{mauve},
	breaklines=true,
	breakatwhitespace=true,
	tabsize=4,
	literate={á}{{\'a}}1 {é}{{\'e}}1 {í}{{\'i}}1 {ó}{{\'o}}1 {ú}{{\'u}}1 {Á}{{\'A}}1 {É}{{\'E}}1 {Í}{{\'I}}1 {Ó}{{\'O}}1 {Ú}{{\'U}}1
}

\begin{document}
	
	\renewcommand{\baselinestretch}{1}
	
	\graphicspath{{./Imagenes/}}
	
	\title{Metodología y Cálculo del Bel de Riesgo para la Reserva de Riesgos en Curso de una Compañía Aseguradora}
	\author{Christopher Gómez Yáñez}
	\institute{Facultad de Ciencias}
	\degree{Actuario}
	\supervisor{Mtro. Alfonso Parrao Guzmán}
	\city{Ciudad Universitaria, CD. MX.}
	\degreemonth{Abril}
	\degreeyear{2020}
	\maketitle
	
	\newpage
	$\ $
	\thispagestyle{empty} % para que no se numere esta pagina
	
	\begin{changemargin}{1cm}{0cm}{1cm}
		
		\vspace{30cm} 
		\begin{center}
			\textit{\textbf{\Large JURADO ASIGNADO}}
		\end{center}
		\vspace{1cm}
		 
		\begin{description}
			\item[]\textbf{Datos del alumno:}\\
			Gómez Yáñez, Christopher\\
			\textit{Número de cuenta:} 307228305\\
			\textit{Carrera:} Actuaria\\
			\textit{Correo:} chris.gomez@ciencias.unam.mx\\
			\textit{Teléfono:} 5532358625
			
			\item[]\textbf{Presidente:}\\
			XYXYXYXYXYXY, Darío\\
			\textit{Correo:} XYXYXYXYXYXY.unam.mx\\
			\textit{Institución de adscripción:} XYXYXYXYXYXY, UNAM
			\item[]\textbf{Vocal:}\\
			XYXYXYXYXYXY, Tonatiuh\\
			\textit{Correo:} XYXYXYXYXYXY.cinvestav.mx \\
			\textit{Institución de adscripción:}XYXYXYXYXYXY, CINVESTAV
			\item[]\textbf{Secretario:}\\
			XYXYXYXYXYXY, Miguel\\
			\textit{Correo:} XYXYXYXYXYXY.unam.mx\\
			\textit{Institución de adscripción:} XYXYXYXYXYXY, UNAM
			\item[]\textbf{1\textsuperscript{er} Suplente:}\\
			XYXYXYXYXYXY, Emilio\\
			\textit{Correo:} XYXYXYXYXYXY.unam.mx\\
			\textit{Institución de adscripción:}XYXYXYXYXYXY, UNAM
			\item[]\textbf{2\textsuperscript{do} Suplente:}\\
			XYXYXYXYXYXY, Juan Carlos\\
			\textit{Correo:} XYXYXYXYXYXY.unam.mx\\
			\textit{Institución de adscripción:}XYXYXYXYXYXY, UNAM
		\end{description}
		\thispagestyle{empty}
	\end{changemargin}
	
	
	%\ChNameVar{\bfseries\Large\sf} \ChNumVar{\Huge} \ChTitleVar{\bfseries\Large\rm}
	%\ChRuleWidth{1pt} \ChNameUpperCase \ChTitleUpperCase
	\ChNumVar{\fontsize{50}{50}\usefont{T1}{ptm}{m}{sl}\selectfont}
	\ChTitleVar{\raggedright\Large\sffamily\bfseries}
	
	\evensidemargin 0in 
	\oddsidemargin 0.6in
	
	\newpage{\ } 
	\thispagestyle{empty}
	
	\begin{dedication}
		{\Large{\sffamily{El continuo esfuerzo, no la fortaleza o inteligencia, es la clave para desbloquear nuestro potencial.}}}\\
		\vspace{0.5cm}
		{\normalsize{\bfseries{Winston S. Churchill}}}
	\end{dedication}
	
	\newpage
	$\ $
	\thispagestyle{empty} % para que no se numere esta pagina
	
	\begin{acknowledgements}
	\pagenumbering{Roman}
	 
	A XXXXXXXX
	\\

	A XXXXXXXXXXX.
	\\
	
	A XXXXXXXXXXXXXXXXXXXXXX.
	\\
	
	AXXXXXXXXXXXXXXXXXXXXX.
	\\
	
	A XXXXXXXXXXXXXXXXXXXXXX.
	\\
	
	A XXXXXXXXxxXXXXXXXXXXx.
	\\
	
	A XXXXXXXxxxxxxxxxxxxx.   
	\\
	
	A XXXXXXxxxxxxxxXXXXXXXXXXXXXX
	\\
	
	A XXXXXXXXXXXXXXXXx
	\\
	
	A XXXXXXXXXXXXXXXXXXXXXXXXx
	\\
	
		
	\end{acknowledgements}
	
	
	\tableofcontents
	
	
	\addtolength{\headheight}{\baselineskip}
	\fancyhead[LE]{\scshape\thepage\hspace{1cm}\footnotesize\nouppercase{Universidad Nacional Autónoma de México}}
	\fancyhead[RO]{\scshape\footnotesize\nouppercase{Facultad de Ciencias}\hspace{1cm}\normalsize\thepage}
	\pagestyle{fancy}
	\cleardoublepage
		
	\begin{preface}
	\pagenumbering{arabic}
	\addcontentsline{toc}{chapter}{\numberline{}Prefacio}%	
	\doublespacing
	
	En este trabajo se tiene como objetivo describir brevemente las labores de un Actuario de Reservas dentro de una Compañía Aseguradora, los cuales incluyen pasos escenciales para poder realizar la valuación de la Reserva de Riesgos en Curso, la cual es la función principal del Actuario. También se mencionaran los reportes que posteriormente harán uso de los resultados de dicha valuacion.
		
	\end{preface}
	
	%--------------------------------------------------------------------------------------------------------------- %
	
	\fancypagestyle{plain}{
		\fancyhead[L]{}
		\fancyhead[C]{}
		\fancyhead[R]{}
		
		\fancyfoot[L]{}
		\fancyfoot[C]{\thepage}
		\fancyfoot[R]{}
		\renewcommand{\headrulewidth}{0pt}
		\renewcommand{\footrulewidth}{0pt}
	}
	
	\fancyhead[LE]{\scshape\thepage\hspace{1cm}\footnotesize\nouppercase{\leftmark}}
	\fancyhead[RO]{\scshape\footnotesize\nouppercase{Sección \rightmark\hspace{1cm}}\normalsize\thepage}
	\fancyhead[LO]{}
	\fancyhead[RE]{}
	\pagestyle{fancy}
	\cleardoublepage


	\chapter{Conceptos Básicos}

Antes que nada hay que familiarizarnos con ciertos conceptos, herramientas y metodologías necesarias para el proceso así como las leyes y organismos gubernamentales que dictan las condiciones bajo las cuales debe realizarse el cálculo.

	\section{El Seguro}
	El seguro es un medio para la protección individuos frente a las consecuencias de riesgos y se basa en transferir dichos riesgos a la institución de seguros, la cual se encargará de indemnizar todo o parte del perjuicio que se produzca por la ocurrencia de un evento previsto\footnote{Ver \citet{ASeguro}, Fundación Mapfre, El Seguro}.
	
	\section{Instituciones de Seguros}

	Una Institución de Seguros es una sociedad anónima autorizada para organizarse y operar conforme a la Ley de Instituciones de Seguros y de Fianzas (LISF), como institución de seguros\footnote{Ver \citet{BAseguradora}, Artículo 2 Sección XVI}.

 
	
	\section{Contrato de Seguro}
	
	El Contrato de Seguro es aquel con que la Empresa Aseguradora se obliga, mediante el pago de una prima, a resarcir un daño o a pagar una suma de dinero al verificarse la eventualidad prevista en el contrato\footnote{Ver \citet{CContrato}, Artículo 1}.
	
		\section{Comisión Nacional de Seguros y Fianzas}
	
	La Comisión Nacional de Seguros y Fianzas es un Órgano Desconcentrado de la Secretaría de Hacienda y Crédito Público, encargada de supervisar que la operación de los sectores asegurador y afianzador se apegue al marco normativo, preservando la solvencia y estabilidad financiera de las instituciones de Seguros y Fianzas, para garantizar los intereses del público usuario, así como promover el sano desarrollo de estos sectores con el propósito de extender la cobertura de sus servicios a la mayor parte posible de la población\footnote{Ver \citet{EComision}}. 
	
	\section{Secretaría de Hacienda y Credito Publico}
	
	La Secretaría de Hacienda y Crédito Público es la dependencia del Poder Ejecutivo Federal que tiene como misión proponer, dirigir y controlar la política económica del Gobierno Federal en materia financiera, fiscal, de gasto, de ingresos y deuda pública, con el propósito de consolidar un país con crecimiento económico de calidad\footnote{Ver \citet{NSHCP}}.
	
	
	\section{Asociación Mexicana de Instituciones de Seguros}
	
	Organismo gremial que representa al interés general de las compañías aseguradoras, promoviendo el desarrollo sano y sustentable del seguro a través de las mejores prácticas. Su principal objetivo es promover el desarrollo de la industria aseguradora, representar sus  intereses ante autoridades del sector público, privado y social, así como proporcionar apoyo técnico a sus asociados\footnote{Ver \citet{FAmis}}. 
	
	\section{Reserva Técnica}
Son las reservas ligadas directamente con los riesgos que se encuentran en curso, incluyendo obligaciones pendientes, provisiones para contingencias y fondos catastróficos\footnote{Ver \citet{RResAssal}}.
De acuerdo a la LISF, las reservas técnicas deberán constituirse y valuarse de forma prudente, confiable y objetiva, en relación con todas las obligaciones de seguro que las Instituciones de Seguros asuman frente a los asegurados y beneficiarios del contrato de seguro, los gastos de administración, si como los gastos de adquisición que, en su caso, asuman con relación a los mismos. Para la constitución se deben utilizar métodos actuariales con base en la aplicación de los estándares de práctica actuarial, considerando la información disponible en los mercados financieros, así como la que generalmente se encuentra disponible sobre riesgos técnicos de seguros\footnote{Ver \citet{DReservasTec}, Artículo 217 y 218}. 


	\section{Reserva de Riesgos en Curso (RRC)}
Constituye una parte de la Reserva Técnica. La Reserva de Riesgos en Curso se puede definir técnicamente como la parte de la prima que debe ser utilizada para el cumplimiento de las obligaciones futuras por concepto de reclamaciones, a lo que también se le llama 
“Prima no Devengada”. Esta definición se traduce a una serie de cálculos actuariales que pueden resultar de mayor o menor complejidad dependiendo del tipo de seguro, por lo que para hacer una mejor explicación es importante hacer la distinción de los diferentes tipos de seguros y referirnos a los métodos de constitución de reservas correspondientes a cada uno de ellos\footnote{Ver \citet{RResAssal}}. El propósito de la RRC es cubrir el valor esperado de las obligaciones futuras derivadas del pago de siniestros, beneficios, valores garantizados, dividendos, gastos de adquisición y administración, así como cualquier otra obligación futura derivada de los contratos del seguro.
La constitución, incremento, valuación y registro de la reserva de riesgos en curso, deberá efectuarse mediante la estimación de obligaciones que se realice empleando los métodos actuariales que, según corresponda, las Instituciones de Seguros y Sociedades Mutualistas registren para tales efectos ante la Comisión y apegándose a los principios y lineamientos establecidos en las Disposiciones de la CUSF. Así mismo, el monto de la reserva de riesgos en curso será igual a la suma de la mejor estimación y de un margen de riesgo, los cuales deberán calcularse por separado y en términos de lo previsto en la CUSF\footnote{Ver \citet{HCusf}, Apartado 5.1.1. y 5.1.3}.

	\begin{figure}[h!]
	\centering
	\includegraphics[height=7cm]{Imagenes/pikachu.png}
	\caption{Imagen de un pokemon.}
	\label{pokemon}
	\end{figure}
 
	\section{Mejor Estimación de Riesgo (Best Estimate of Loss)}
	
	La mejor estimación será igual al valor esperado de los flujos futuros de obligaciones, entendido como la media ponderada por probabilidad de dichos flujos. Las hipótesis y procedimientos con que se determinen los flujos futuros de obligaciones, con base en los cuales se obtendrá la mejor estimación, deberán ser definidos por la Institución de Seguros en el método propio que registre para el cálculo de la mejor estimación. El cálculo de la mejor estimación se basará en información oportuna, confiable, homogénea y suficiente, así como en hipótesis realistas, y se efectuará empleando métodos actuariales y técnicas estadísticas basados en la aplicación de los estándares de práctica actuarial\footnote{Ver \citet{HCusf}, Apartado 5.1.3}. Ir a figura \ref{pokemon}
	
	\section{Requerimiento de Capital de Solvencia (RCS)}
	
	El Requerimiento de Capital de Solvencia es el requerimiento de los recursos patrimoniales, adicional a la Base de Inversión, con los que la institución debe contar para hacer frente a las obligaciones con los asegurados, derivados de desviaciones no esperadas relacionadas con los riesgos técnicos, de reaseguro, financieros y operativos. Las inversiones que respaldan este requerimiento deben encontrarse en condiciones adecuadas de seguridad y liquidez conforme a la regulación aplicable\footnote{Ver \citet{NRCS}}.
	
	
		\section{Lenguaje R}
	R es un lenguaje y ambiente computacinal enfocado en estadistica y graficas. Se trata de una herramienta informática sumamente potente para realizar distintos cálculos científicos, numéricos y estadísticos, así como para crear gráficas y figuras de gran calidad. R es un programa gratuito, relativamente fácil de operar y cuenta con una gran comunidad de internet que contribuye a resolver dudas y problemas, sin costo alguno\footnote{Ver \citet{KR}}.
	
	\chapter{Marco Regulatorio}\label{tcyedb}
	
	Para calcular la Reserva de Riesgos en Curso deben de seguirse los lineamientos de los organismos reguladores del gobierno, los cuales incluyen Leyes, Circulares y Anexos que estan sujetos a cambios segun lo determinen dichos organizmos.
	
	\section{Circular Única de Seguros y Fianzas}

	
	Cuerpo normativo que contiene las disposiciones derivadas de la Ley de Instituciones de Seguros y de Fianzas, que dan operatividad a sus preceptos y sistematizan su integración, homologando la terminología utilizada, a fin de brindar con ello certeza jurídica en cuanto al marco normativo al que las instituciones y sociedades mutualistas de seguros, instituciones de fianzas y demás personas y entidades sujetas a la inspección y vigilancia de la Comisión Nacional de Seguros y Fianzas deberán sujetarse en el desarrollo de sus operaciones\footnote{Ver \citet{IDefCusf}} .El 19 de diciembre de 2014, se publicó en el Diario Oficial de la Federación la Circular Única de Seguros y Fianzas (CUSF). Esta circular instrumenta y da operatividad a la nueva Ley de Instituciones de Seguros y de Fianzas (LISF) promulgada el 4 de abril de 2013 y en vigor desde el 4 de abril de 2015. 
	
	\section{Ley de Instituciones de Seguros y de Fianzas}

	Ley que tiene por objeto la organización, operación y funcionamiento de las Instituciones de Seguros, Instituciones de Fianzas y Sociedades Mutualistas de Seguros; las actividades y operaciones que las mismas podrán realizar, así como las de los agentes de seguros y de fianzas, y demás participantes en las actividades aseguradora y afianzadora en protección de los intereses del público usuario de estos servicios financieros\footnote{Ver \citet{GLisf}, Artículo 1}. De acuerdo con La Ley de Instituciones de Seguros y de Fianzas se establece que las Instituciones de Seguros deberán registrar ante la Comisión Nacional de Seguros y Fianzas los métodos actuariales con base en sus estimaciones para la Reserva de Riesgos en Curso\footnote{Ver \citet{FAmis}, Artículo 219}, de conformidad con las disposiciones de carácter general que al efecto emita, mismas que se dieron a conocer a través de la Circular Única de Seguros y Fianzas publicada en el Diario Oficial de la Federación\footnote{Ver \citet{HCusf}}.
	 
	
	
\chapter{Funciones Generales}\label{cap:Introducción}

Las funciones de un Actuario de Reservas se distinguen por tres actividades principales: validaciones de calidad de la información, uso de dicha información para el cálculo de cifras financieras y, finalmente, uso de dichas cifras para el llenado de reportes para el regulador. A continuación se mencionan brevemente el desarrollo de estas actividades.

\begin{singlespace}
	\begin{enumerate}
		\item \textbf{Actividades Mensuales}
		\begin{enumerate}
			\item \textit{Procesos del Cierre de Mes} \begin{enumerate}
				\item Revisión de bases de información de Cierre:
				\\ \-\hspace{0.5cm} El Actuario debe asegurarse de que la información correspondiente a los Asegurados y Siniestros registrados en las bases de datos de la Compañía sea veraz, oportuna, integra y consistente, con el fin de utilizarla en procesos consiguientes, como son reportes y especialmente para el cálculo de la Reserva Técnica de la Compañía.\\
			\end{enumerate}
			\item \textit{ Cálculo de Reserva Técnica}  	\begin{enumerate}
				\item Verificar Factores de Mercado publicados en la CUSF para calculo de Reservas mediante Metodo Estatutario:\\ \-\hspace{0.5cm}
				Para el Cálculo de Reservas de grupos homogeneos cuya información no es suficiente se debe emplear el metodo Estatutario aprovado por la CNSF, para el cual son necesarios Factores que se obtienen mediante el uso de la Información de Mercado de Seguros. Ya que el Mercado es cambiante es responsabilidad del Actuario asegurarse que se estan empleando los factores más actualizados al momento del cálculo.
				\item Extracción de información de las Bases de Datos validadas:\\ \-\hspace{0.5cm}
				Para calcular la Reserva es necesario contar con las cifras correspondientes a Asegurados en Vigor y Siniestros Autorizados, para la obtención de la Reserva de Gastos de Administración y la Reserva de Riesgos en Curso, respectivamente.				
				\item Simulación Estocástica en R:\\ \-\hspace{0.5cm}
				Una vez que se cuenta con la información de Siniestralidad Autorizada se generara a partir de élla las Matrices de Siniestralidad sobre las cuales se aplicara el método expuesto en este trabajo, en específico las $n$ simulaciones.
				\item Obtención de los resultados del Cálculo de la Reserva:\\ \-\hspace{0.5cm}
				Se realizara un resumen con los resultados del Cálculo en una plantilla de Excel, el cual incorporara el Cálculo de la primera Reserva, así como la conclusión de las $n$ simulaciones posteriores.
				\item Envío de Información a Áreas Técnicas:\\ \-\hspace{0.5cm}
				Se enviaran los resultados de Reserva Técnica al resto de Áreas Técnicas para su registro en los reportes internos y su uso posterior en entregas de información al Organizmo Regulador\\
				
			\end{enumerate}
			\item \textit{Obtención de información para el Cálculo del Requerimiento de Capital de Solvencia (RCS)} 
			\begin{enumerate} 
				\item Obtención de Factóres de Devengamiento de la CUSF:\\ \-\hspace{0.5cm}
				Para el Cálculo de Insúmos de Suscripción es necesario contar con los factóres de Devengamiento publicados en la última actualización del \textit{Manual de datos para el cálculo del RCS de índices de siniestralidad del mejor estimador}.
				\item Obtención de Insúmos de Suscripción:\\ \-\hspace{0.5cm}
				Una vez obtenidos los Factóres de Devengamiento actualizados se calcúlan los Insumos de Suscripción mediante el uso de la Reserva Técnica calculada y los Asegurados en Vigor al Cierre de mes.
				\item Envío de los Insumos de Suscripción al Área Responsable del Cálculo del nuevo RCS:\\ \-\hspace{0.5cm} Una vez que se calcúlan los Insumos de Suscripción, estos deben enviarse al Area Responsable del Cálculo del RCS para que esta los ingrese al validador en linea de la CNSF y posteriormente envíe el nuevo RCS en el reporte correspondiente.\\ 
			\end{enumerate}
		\end{enumerate}
		\item \textbf{Actividades Trimestrales}
		\begin{enumerate}
			\item \textit{Entrega de Reportes Trimestrales a la CNSF} \\ \-\hspace{0.5cm}
			La CNSF hace del conocimiento de las Instituciones de Seguros y Fianzas la responsabilidad de entregar reportes trimestrales con información que sustente la correcta operación de dichas Compañías.
			\begin{enumerate}
				\item {Ingresar a la pagina oficial de la CNSF para obtener las fechas límite de entrega de los diferentes reportes} \\ \-\hspace{0.5cm}
				La CNSF actualiza periodicamente el listado de reportes que se deben entregar, así como las fechas antes de las cuales debe hacerse el envío. Es responsabilidad del actuario estar informado de dichas actualizaciones para enviar los reportes en tiempo y forma.
				\item {Llenado y entrega del REPORTE REGULATORIO SOBRE RESERVAS TÉCNICAS Trimestral (RR-3 REVAL)} \\ \-\hspace{0.5cm}
				Siguiendo el Manual oficial de la CNSF debe ser llenado el reporte RR3 REVAL, el cuál esta consituido por la información del cálculo de Reservas Técnicas (Comparte información con el REPORTE REGULATORIO SOBRE ESTADOS FINANCIEROS RR7 por lo que se debe validar que no existan inconsistencias entre los reportes). Posteriormente, la información debe ser depositada en el Sistema de Entrega de Información Vía Electronica (SEIVE).
				\item {Llenado y entrega del Reporte de Información Estadística por Operación, Ramo o Seguro (RR8 ORS)} \\ \-\hspace{0.5cm}
				Siguiendo el Manual oficial de la CNSF debe ser llenado el reporte RR8 ORS, el cuál esta consituido por la información estadística de los Asegurados en la cartera de Compañía. Posteriormente, la información debe ser depositada en el Sistema de Entrega de Información Vía Electronica (SEIVE).
				\item {Llenado y entrega del Reporte de Información del Comportamiento por Operación y Ramo (RR8 COR)} \\ \-\hspace{0.5cm}
				Siguiendo el Manual oficial de la CNSF debe ser llenado el reporte RR8 COR, el cuál esta consituido por el número de Pólizas y Asegurados y Monto de Prima por Subramo en la Cartera de la Compañía. Posteriormente, la información debe ser depositada en el Sistema de Entrega de Información Vía Electronica (SEIVE).
			\end{enumerate}
			\item \textit{Entrega de Reportes Trimestrales a la AMIS} \\ \-\hspace{0.5cm}
			Con el fin de realizar análisis estadisticos del mercado en beneficio de las instituciones de Seguros y Fianzas, la AMIS solicita a las Compañías que le envíen los reportes Trimestrales RR mediante un Repositorio en Linea. Es responsabilidad del Actuario hacer el envío Trimestral de la información una vez se hayan entregado de manera exitosa los Reportes a la CNSF.
			\item \textit{Entrega de información al Auditor Externo} \\ \-\hspace{0.5cm}
			Con el propósito de dar cumplimiento a lo establecido por la CNSF, las Compañías estan obligadas a contratar un Auditor Externo cuya función será buscar inconsistencias en su información interna para su pronta y oportuna corrección. Una vez corregida la información el Auditor Externo ratificará ante la CNSF que esta es confiable. Es responsabilidad del Actuario compartir con el Auditor Externo la información convenida en el Programa de Auditoría para posteriormente, en conjunto, resolver dudas y acordar soluciones. 
		\end{enumerate}
		\item \textbf{Actividades Anuales}
		\begin{enumerate}
			\item \textit{Asistencia a los foros organizados por la AMIS para revisar nuevos reportes o modificaciones a reportes anuales que solicita la CNSF, aclarar dudas y acordar propuestas para presentar a la CNSF.} \\ 
			\item \textit{Entrega de Reportes Trimestrales a la CNSF} \\ \-\hspace{0.5cm}
			La CNSF hace del conocimiento de las Instituciones de Seguros y Fianzas la responsabilidad de entregar reportes anuales con información que sustente la correcta operación de dichas Compañías.
			\begin{enumerate}
				\item {Ingresar a la pagina oficial de la CNSF para obtener las fechas límite de entrega de los diferentes reportes} \\ \-\hspace{0.5cm}
				La CNSF actualiza periodicamente el listado de reportes que se deben entregar, así como las fechas antes de las cuales debe hacerse el envío. Es responsabilidad del actuario estar informado de dichas actualizaciones para enviar los reportes en tiempo y forma.
				\item {Llenado y entrega del REPORTE REGULATORIO SOBRE RESERVAS TÉCNICAS Anual (RR-3 DETAN)} \\ \-\hspace{0.5cm}
				Siguiendo el Manual oficial de la CNSF debe ser llenado el reporte RR3 DETAN, el cuál esta consituido por la información del cálculo de Reservas Técnicas de cada asegurado en la cartera de la Compañía (Comparte información con el REPORTE REGULATORIO SOBRE ESTADOS FINANCIEROS RR7 por lo que se debe validar que no existan inconsistencias entre los reportes). Posteriormente, la información debe ser depositada en el Sistema de Entrega de Información Vía Electronica (SEIVE).
				\item {Llenado y entrega del Reporte SISTEMA ESTADÍSTICO DEL SECTOR ASEGURADOR (RR8 SESA)} \\ \-\hspace{0.5cm}
				Siguiendo el Manual oficial de la CNSF debe ser llenado el reporte RR8 SESA para cada ramo que la Compañía maneje en su cartera, el cuál esta consituido por la información estadística de cada Asegurado en la cartera de la Compañía. Posteriormente, la información debe ser depositada en el Sistema de Entrega de Información Vía Electronica (SEIVE).
				\item {Llenado y entrega del Reporte FORMAS ESTADÍSTICAS DE LOS SEGUROS (RR8 FES)} \\ \-\hspace{0.5cm}
				Siguiendo el Manual oficial de la CNSF debe ser llenado el reporte RR8 FES para cada ramo que la Compañía maneje en su cartera, el cuál esta consituido por la información estadística por estado de cada Asegurado en la cartera de la Compañía. Posteriormente, la información debe ser depositada en el Sistema de Entrega de Información Vía Electronica (SEIVE).
					\item {Llenado y entrega del Reporte Regulatorio sobre Reservas Técnicas (RR3 RRCBD)} \\ \-\hspace{0.5cm}
				Siguiendo el Manual oficial de la CNSF debe ser llenado el reporte RR3 RRCBD. El Backtesting es un procedimiento técnico que consiste en validar la idoneidad, precisión y validez de un modelo diseñado para hacer estimaciones, con base en la comparación de las estimaciones hechas por el modelo actuarial respecto de los valores reales observados. Posteriormente, la información debe ser depositada en el Sistema de Entrega de Información Vía Electronica (SEIVE).
			\end{enumerate}
		
			\item \textit{Entrega regulatoria de la carta de aceptación del Actuario Independiente firmada en conjunto del Contrato de Servicios del Actuario Independiente} \\ \-\hspace{0.5cm}
			Como parte de la regulación, es necesario hacer el envío a la CNSF de la carta de aceptación del Actuario que fungira como Auditor Externo de la Compañía con las certificaciones correspondientes. De igual forma debe incluirse el contrato firmado en el que se define el periodo, terminos y condiciones para que el Actuario Independiente cumpla con sus funciones como Auditor Externo. La información debe ser depositada en el Sistema de Entrega de Información Vía Electronica (SEIVE).
			\item \textit{Entrega regulatoria del Programa de Actividades del Actuario Independiente} \\ \-\hspace{0.5cm}
			Como parte de la regulación, es necesario hacer el envío a la CNSF del Programa de Actividades del Actuario Independiente, el cual debe incluir la información que se auditará, así como los periodos de entrega de información por parte de la Compañía y las fechas de entrega de comentarios por parte del Auditor. 
		\end{enumerate}
		\item \textbf{Otras actividades}
		\begin{enumerate}
			\item \textit{Automatización y Documentación de Procesos} \\ \-\hspace{0.5cm}
			Debido a la cantidad y frecuencia de reportes regulatorios, es responsabilidad del Actuario buscar nuevas y mejores formas de procesar y extraer la información que se requiere en la operación diaria, así como documentar dichos procesos para facilitar su uso.
			\item \textit{Creación de Manuales de Capacitación} \\ \-\hspace{0.5cm}
			El Actuario debe constantemente crear y/o actualizar manuales que detallen sus funciones para la facil asimilacion de dichas funciones por nuevos miembros de equipo, al igual que los Auditores Internos de la Compañía para así asegurar la correcta operación y la calidad de la información.
		\end{enumerate}
	\end{enumerate}
\end{singlespace}




	\chapter{Metodología empleada en el Cálculo de la Reserva de Riesgos en Curso}\label{tcyedb}
	
	La Metodología son los fundamentos y técnicas que se aplican durante el proceso del Cálculo de la Reserva de Riesgos en Curso para alcanzar un resultado válido. Funciona como el soporte conceptual ante el Regulador que justifica la manera en que aplicamos los procedimientos en el Cálculo.
	
	\section{Método Chain-Ladder}

	Método que se utiliza comunmente en las reservas de no-vida. Útiliza un factor para "suavizar" los datos y con base en estos, realizar interpolaciones para estimar los siniestros agregados para cada año de ocurrencia y posteriormente la reserva correspondiente. El supuesto básico de este método es que las columnas en el triángulo de desarrollo son proporcionales, es decir que, independientemente del año de origen, cada periodo de desarrollo se reporta una proporción constante de siniestros con respecto al total. La sustentación del supuesto depende en buena medida, tanto del tipo de negocio que se trate, como de la homogeneidad y tamaño de la cartera. En particular, en negocios como vida individual, gastos médicos, responsabilidad civil, etc., la evolución del reporte de los siniestros es estacional\footnote{Ver \citet{LChainLadder}, 1.2.1 Método Chain-Ladder}.  

	La estimación de las obligaciones se hace con base en los siniestros observados y su desfase respecto a la entrada en vigor de cada obligación, usando el metodo $bootstrap$.
	
	\section{Bootstrapping}
	El método desarrollado por Bradley Efron\footnote{Ver \citet{OEfron}}. Es un método de muestreo computacionalmente intensivo con el que se busca aproximar la distribución muestral de alguna variable aleatoria que se basa en los datos observados\footnote{Ver \citet{MBootstrap}, p.11}.

	El método de $Bootstrap$ es un método de muestreo con el que se busca aproximar la distribución muestral de alguna variable aleatoria que tiene como base los datos observados.

	Teniendo una muestra de datos $x_{1},x_{2},x_{3},...,x_{n}$, donde los $x_{i}$ son independientes y provienen de una distribución desconocida $F$, donde además se presume que dicha muestra es una representación significativa de la población de donde proviene. Se tiene además una variable aleatoria $R(X,F)$ que depende de $X$ y de la función desconocida $F$. Entonces se puede realizar una muestra aleatoria de tamaño $n$ con reemplazo de la muestra de datos, $x_{1}^{*},x_{2}^{*},x_{3}^{*},...,x_{n}^{*}$ y a partir de esa muestra se puede calcular una observación de la variable aleatoria $R*(X*,P*)$, donde $F*$ es la distribución de probabilidad de la muestra, que se construyó de tipo uniforme. Finalmente, se realizan más muestras y se calculan más valores de $R*$ para poder estimar la distribución $R(X,F)$.

	La utilidad técnica de bootstrapping es que permite aproximar la distribución de alguna estadística de los datos de una forma fácil y rápida. Adicionalmente, no es necesario hacer una estimación paramétrica ni supuestos acerca de la distribución de los datos.
	
	


		
	\chapter{Cálculo del Bel de Riesgos en Curso}

 La valuación y constitución de la reserva de riesgos en curso esta compuesta por el Mejor Estimador de Riesgo, la Reserva de Gasto y el Margen de Riesgo. El presente trabajo se enfóca en la obtención del Mejor Estimador, presentado a continuación.
 
 El proceso descrito debe efectuarse para un grupo homogéneo definido, correspondiente a un cierto subramo y tipo de seguro que la Compañia en cuestión tenga en su cartera y puede aplicarse en la practica para grupos homogéneos correspondientes a los ramos de No Vida.

	\section{Aplicación del Método Chain-Ladder}
	
	El cálculo del Bel de Riesgo, siguiendo el método $Chain-Ladder$, implica un análisis de las obligaciones futuras para los riesgos en curso con base en los siniestros que actualmente han sido reportados. Para ello se necesita la construcción de una matriz de desarrollo de siniestros  de dimensiones $(k$ x $s)$, en la cual los siniestros se distribuyen por el trimestre en que se reporto cada uno de los procedimientos ocurridos respecto al inicio de vigencia de la póliza. La matriz queda como se muestra en la tabla ~\ref{matrix1}.

	\vspace{0.1cm}
	
	\begin{itemize}
		\setlength\itemsep{-0.5em}
		\item ${X}_{i,j}$ es el monto de siniestros de las pólizas con inicio de vigencia en el trimestre $i$ que fue reportado $j$ trimestres posteriores al inicio de vigencia.
		\item ${k}_{}$ es el número de trimestres máximo observado en la experiencia de siniestros.
		\item ${s}_{}$ es el número de trimestres de experiencia de inicio de vigencia.
		\item $i$ es el trimestre de inicio de vigencia de la póliza, $i\in \left\{1,2,3,\dots ,12\right\}$
		\item $j$ es el trimestre en que se reportó el siniestro,  $j\in \left\{1,2,3,\dots\right\}$
	\end{itemize}
	
	\vspace{1cm}

	\begin{table}[ht]
	\centering
	\begin{tabularx}{\linewidth}{c|ccccccccc}
		\multirow{2}{4cm}{Trimestre de inicio de vigencia de la póliza} & \multicolumn{9}{c}{Trimestre en que se reportó el procedimiento} \\
			& 0  & 1 & 2 & $ \dots $ & j & $\dots $ & k-2 & k-1 &  k \\
		\midrule
		1      &  $X_{1,0}^{}$ & $X_{1,1}^{}$ & $X_{1,2}^{}$ & $ \dots $ & $X_{1,j}^{}$ & $ \dots $ & $X_{1,k-2}^{}$ & $X_{1,k-1}^{}$ & $X_{1,k}^{}$ \\
		2      &  $X_{2,0}^{}$ & $X_{2,1}^{}$ & $X_{2,2}^{}$ & $ \dots $ & $X_{2,j}^{}$ & $ \dots $ & $X_{2,k-2}^{}$ & $X_{2,k-1}^{}$ & \\
		3      &  $X_{3,0}^{}$ & $X_{3,1}^{}$ & $X_{3,2}^{}$ & $ \dots $ & $X_{3,j}^{}$ & $ \dots $ & $X_{3,k-2}^{}$ & & \\
		4      &  $X_{4,0}^{}$ & $X_{4,1}^{}$ & $X_{4,2}^{}$ & $ \dots $ & $X_{4,j}^{}$ & $ \dots $ & & & \\
		:      & & & & & & & & &\\
		i      &  $X_{i,0}^{}$ & $X_{i,1}^{}$ & $X_{i,2}^{}$ & $ \dots $ & $X_{i,j}^{}$ & & & & \\
		:      & & & & & & & & & \\
		s-2      &  $X_{s-2,0}^{}$ & $X_{s-2,1}^{}$ & $X_{s-2,2}^{}$ & & & & & & \\
		s-1      &  $X_{s-1,0}^{}$ & $X_{s-1,1}^{}$ & & & & & & & \\
		s      &  $X_{s,0}^{}$ & & & & & & & & \\
	\end{tabularx}
\caption{Matriz de Siniestros}
\label{matrix1}
	\end{table}

	\vspace{1cm}

	En este punto se puede usar la esta matriz de siniestros para generar una matriz de siniestros acumulados (tabla \ref{matrix2}) en la cual se define ${Y}_{i,j}$ como el monto de siniestros de la póliza con inicio de vigencia en el trimestre $i$ reportados hasta el trimestre $j$, como se ve en \citet{PAlvColl} (ver eq. \ref{eq1}).
	
	\begin{equation}
	{y}_{i,j}=\sum _{m=0}^{j}{X}_{i,m}
		\label{eq1}
	\end{equation}

	Donde
	
	\begin{itemize}
	\setlength\itemsep{-0.5em}
		\item ${X}_{i,m}$ es el monto de siniestros de las pólizas con inicio de vigencia en el trimestre $i$ que fue reportado $m$ trimestres posteriores al inicio de vigencia.
		\item $i$ es el trimestre de inicio de vigencia de la póliza, $i\in \left\{1,2,3,\dots ,12\right\}$
		\item $j$ es el trimestre en que se reportó el siniestro,  $j\in \left\{1,2,3,\dots\right\}$
		\item $m$ es el trimestre de acumulación, $m\in \left\{0,1,2,\dots ,j\right\}$, $m\le j\le k$
	\end{itemize} 
	
	\begin{table}[ht]
	\centering
		
	\begin{tabularx}{\linewidth}{c|cccccccccc}
		\multirow{2}{4cm}{Trimestre de inicio de vigencia de la póliza}&\multicolumn{9}{c}{ Trimestre en que se reportó el procedimiento} \\
		& 0  & 1 & 2 & $ \dots $ & j & $\dots $ & k-2 & k-1 &  k \\
		\midrule
		1      &  $Y_{1,0}^{}$ & $Y_{1,1}^{}$ & $Y_{1,2}^{}$ & $ \dots $ & $Y_{1,j}^{}$ & $ \dots $ & $Y_{1,k-2}^{}$ & $Y_{1,k-1}^{}$ & $Y_{1,k}^{}$ \\
		2      &  $Y_{2,0}^{}$ & $Y_{2,1}^{}$ & $Y_{2,2}^{}$ & $ \dots $ & $Y_{2,j}^{}$ & $ \dots $ & $Y_{2,k-2}^{}$ & $Y_{2,k-1}^{}$ & \\
		3      &  $Y_{3,0}^{}$ & $Y_{3,1}^{}$ & $Y_{3,2}^{}$ & $ \dots $ & $Y_{3,j}^{}$ & $ \dots $ & $Y_{3,k-2}^{}$ & & \\
		4      &  $Y_{4,0}^{}$ & $Y_{4,1}^{}$ & $Y_{4,2}^{}$ & $ \dots $ & $Y_{4,j}^{}$ & $ \dots $ & & & \\
		:      & & & & & & & & &\\
		i      &  $Y_{i,0}^{}$ & $Y_{i,1}^{}$ & $Y_{i,2}^{}$ & $ \dots $ & $Y_{i,j}^{}$ & & & & \\
		:      & & & & & & & & & \\
		s-2      &  $Y_{s-2,0}^{}$ & $Y_{s-2,1}^{}$ & $Y_{s-2,2}^{}$ & & & & & & \\
		s-1      &  $Y_{s-1,0}^{}$ & $Y_{s-1,1}^{}$ & & & & & & & \\
		s      &  $Y_{s,0}^{}$ & & & & & & & & \\
	\end{tabularx}
\caption{Matriz de Siniestros Acumulados}
\label{matrix2}
	\end{table} 
	
	
	Usando la esta matriz de siniestros acumulados se obtienen los factores de desarrollo ${f}_{j}$, los cuales indican el incremento dado de un trimestre a otro (ver eq. \ref{eq2}).

	\begin{equation}
	{f}_{j}=\frac{\sum _{i=1}^{s-j}{Y}_{i,j}}{\sum _{i=1}^{s-j}{Y}_{i,j-1}}, \hspace{0.5cm}\text{con}\hspace{0.5cm} 0< j\le k
	\label{eq2}
	\end{equation}

	Donde ${Y}_{i,j}$ representa el monto de siniestros de las pólizas con inicio de vigencia en el trimestre $i$ reportados hasta el trimestre $j$.

	Mediante estos factores de desarrollo, se definen los Siniestros Esperados para la vigencia $i$ (${SE}_{i}$) como la estimación del monto de siniestros que serán reportados para las pólizas con inicio de vigencia en el trimestre $i$ (ver eq. \ref{eq3}).
	
	\begin{equation}
	{SE}_{i}={Y}_{s-i+1,i-1}\cdot\Pi_{j=i}^{k}{f}_{j},\hspace{0.5cm}\text{con}\hspace{0.5cm}0< i\le k
		\label{eq3}
	\end{equation}

	Donde ${Y}_{s-i+1,i-1}$ representa el monto de siniestros de las pólizas con inicio de vigencia en el trimestre $s-i+1$ reportados hasta el trimestre $i-1$, y ${f}_{i}$ es el factor de incremento del trimestre $j$. 

	Se definen las obligaciones futuras iniciales de riesgos en curso (${RRC}^{0}$) como la diferencia entre los Siniestros Estimados y los Siniestros Acumulados observados (ver eq. \ref{eq4}).
	
	\begin{equation}
	{RRC}^{0}=\sum _{i=s-k}^{s}{SE}_{i}-{Y}_{i,s-i}
	\label{eq4}
	\end{equation}

	Donde ${SE}_{i}$ son los siniestros esperados para la vigencia $i$, y ${Y}_{i,s-i}$ es el monto de siniestros de las pólizas con inicio de vigencia en el trimestre $i$ reportados hasta el trimestre $s-i$.
		 
\section{Ajuste de los Triángulos de Siniestralidad}		 
		 
Ya que el método $Chain-Ladder$ puede ser considerado una regresión ponderada, se pueden obtener residuales a partir de un ajuste en los triangulos de siniestralidad usando el método como lo describe \citet{QLowe}, calculando un triángulo de desarrollo acumulado de datos ajustados, partiendo del último dato observado y desplazándose hacia atrás con los factores de desarrollo obtenidos previamente. Se define ${Y}_{i,j}^{*}$ como el monto ajustado de siniestros de las pólizas con inicio de vigencia en el trimestre $i$ reportados hasta el trimestre $j$ (ver eq. \ref{eq5}).
	
		\begin{equation}
	{Y}_{i,j-1}^{*}=\frac{{Y}_{i,j}^{*}}{{f}_{j}}
	\label{eq5}
	\end{equation}

	Con:

	$${Y}_{i,k-i+1}^{*}={Y}_{i,k-i+1}$$	

	Donde ${Y}_{i,j}$ es el monto de siniestros de las pólizas con inicio de vigencia en el trimestre $i$ reportados hasta el trimestre $j$, y ${f}_{j}$ es el factor de incremento del trimestre $j$.

	Con estos montos se obtiene la matriz de siniestros acumulados ajustados partiendo del último dato observado, como se observa en la tabla ~\ref{matrix3}.
		
	\begin{table}[ht]
		\centering

		\begin{tabularx}{\linewidth}{ c |ccccccccc}
			\multirow{2}{4cm}{Trimestre de inicio de vigencia de la póliza}
			& \multicolumn{9}{c}{Trimestre en que se reportó el procedimiento} \\ 
			& 0  & 1 & 2 & $ \dots $ & j & $\dots $ & k-2 & k-1 &  k \\
			\midrule
			1      &  $Y_{1,0}^{*}$ & $Y_{1,1}^{*}$ & $Y_{1,2}^{*}$ & $ \dots $ & $Y_{1,j}^{*}$ & $ \dots $ & $Y_{1,k-2}^{*}$ & $Y_{1,k-1}^{*}$ & $Y_{1,k}^{}$ \\
			2      &  $Y_{2,0}^{*}$ & $Y_{2,1}^{*}$ & $Y_{2,2}^{*}$ & $ \dots $ & $Y_{2,j}^{*}$ & $ \dots $ & $Y_{2,k-2}^{*}$ & $Y_{2,k-1}^{}$ & \\
			3      &  $Y_{3,0}^{*}$ & $Y_{3,1}^{*}$ & $Y_{3,2}^{*}$ & $ \dots $ & $Y_{3,j}^{*}$ & $ \dots $ & $Y_{3,k-2}^{}$ & & \\
			4      &  $Y_{4,0}^{*}$ & $Y_{4,1}^{*}$ & $Y_{4,2}^{*}$ & $ \dots $ & $Y_{4,j}^{*}$ & $ \dots $ & & & \\
			:      & & & & & & & & & \\
			i      &  $Y_{i,0}^{*}$ & $Y_{i,1}^{*}$ & $Y_{i,2}^{*}$ & $ \dots $ & $Y_{i,j}^{}$ & & & & \\
			:      & & & & & & & & & \\
			s-2      &  $Y_{s-2,0}^{*}$ & $Y_{s-2,1}^{*}$ & $Y_{s-2,2}^{}$ & & & & & & \\
			s-1      &  $Y_{s-1,0}^{*}$ & $Y_{s-1,1}^{}$ & & & & & & & \\
			s      &  $Y_{s,0}^{}$ & & & & & & & & \\
		\end{tabularx}
			\caption{Matriz de Siniestros Acumulados Ajustados}
	\label{matrix3}
	\end{table}
	
	Usando esta matriz de siniestros acumulados ajustados se obtiene una nueva matriz de montos $X_{i,j}^{*}$ (tabla ~\ref{matrix4}), que son los siniestros ajustados de las pólizas con inicio de vigencia en el trimestre $i$ que fueron reportados $j$ trimestres posteriores al inicio de vigencia (ver eq. \ref{eq6}).
	
	\begin{equation}
	{X}_{i,j}^{*}={Y}_{i,j}^{*}-\sum _{m=0}^{j-1}{X}_{i,m}^{*}
	\label{eq6}
	\end{equation}	
	 
	Con:
	
	$${X}_{i,0}^{*}={Y}_{i,0}^{*}$$ 
	
	Donde ${Y}_{i,j}^{*}$ es el monto ajustado de siniestros de las pólizas con inicio de vigencia en el trimestre $i$ reportados hasta el trimestre $j$. 
	
	\begin{table}[ht]
		\centering
		\begin{tabularx}{\linewidth}{ c|ccccccccc}
		\multirow{2}{4cm}{Trimestre de inicio de vigencia de la póliza}&  \multicolumn{9}{c}{Trimestre en que se reportó el procedimiento} \\
			& 0  & 1 & 2 & $ \dots $ & j & $\dots $ & k-2 & k-1 &  k \\
			\midrule
			1      &  $X_{1,0}^{*}$ & $X_{1,1}^{*}$ & $X_{1,2}^{*}$ & $ \dots $ & $X_{1,j}^{*}$ & $ \dots $ & $X_{1,k-2}^{*}$ & $X_{1,k-1}^{*}$ & $X_{1,k}^{*}$ \\
			2      &  $X_{2,0}^{*}$ & $X_{2,1}^{*}$ & $X_{2,2}^{*}$ & $ \dots $ & $X_{2,j}^{*}$ & $ \dots $ & $X_{2,k-2}^{*}$ & $X_{2,k-1}^{*}$ & \\
			3      &  $X_{3,0}^{*}$ & $X_{3,1}^{*}$ & $X_{3,2}^{*}$ & $ \dots $ & $X_{3,j}^{*}$ & $ \dots $ & $X_{3,k-2}^{*}$ & & \\
			4      &  $X_{4,0}^{*}$ & $X_{4,1}^{*}$ & $X_{4,2}^{*}$ & $ \dots $ & $X_{4,j}^{*}$ & $ \dots $ & & & \\
			:      & & & & & & & & &\\
			i      &  $X_{i,0}^{*}$ & $X_{i,1}^{*}$ & $X_{i,2}^{*}$ & $ \dots $ & $X_{i,j}^{*}$ & & & & \\
			:      & & & & & & & & & \\
			s-2      &  $X_{s-2,0}^{*}$ & $X_{s-2,1}^{*}$ & $X_{s-2,2}^{*}$ & & & & & & \\
			s-1      &  $X_{s-1,0}^{*}$ & $X_{s-1,1}^{*}$ & & & & & & & \\
			s      &  $X_{s,0}^{*}$ & & & & & & & & \\
		\end{tabularx}
				\caption{Matriz de Siniestros  Ajustados}
				\label{matrix4}
	\end{table}
		
	Con esta información se genera la Matriz de Residuales Brutos $R_{i,j}^{}$ (tabla ~\ref{matrix5}), los cuales estan dados por la diferencia del monto de siniestros observados (el monto original) y el monto de siniestros ajustado  (ver eq. \ref{eq7}).
		
	\begin{equation}
	R_{i,j}^{}= X_{i,j}^{} - X_{i,j}^{*} 
	\label{eq7}
	\end{equation}
	
	Donde 
	
	\begin{itemize}
			\setlength\itemsep{-0.5em}
		\item ${X}_{i,j}^{*}$ = Monto de siniestros ajustados de las pólizas con inicio de vigencia en el trimestre $i$ que fueron reportados $j$ trimestres posteriores al inicio de vigencia.
		\item ${X}_{i,j}$ = Monto de siniestros de las pólizas con inicio de vigencia en el trimestre $i$ que fueron reportados $j$ trimestres posteriores al inicio de vigencia.		
		\end{itemize}
		

	\begin{table}[ht]
		\centering
		\begin{tabularx}{\linewidth}{ c|ccccccccc}
			\multirow{2}{4cm}{Trimestre de inicio de vigencia de la póliza}
			& \multicolumn{9}{c}{Trimestre en que se reportó el procedimiento} \\
			& 0  & 1 & 2 & $ \dots $ & j & $\dots $ & k-2 & k-1 &  k \\
			\midrule
			1      &  $R_{1,0}^{ }$ & $R_{1,1}^{ }$ & $R_{1,2}^{ }$ & $ \dots $ & $R_{1,j}^{ }$ & $ \dots $ & $R_{1,k-2}^{ }$ & $R_{1,k-1}^{ }$ & $R_{1,k}^{ }$ \\
			2      &  $R_{2,0}^{ }$ & $R_{2,1}^{ }$ & $R_{2,2}^{ }$ & $ \dots $ & $R_{2,j}^{ }$ & $ \dots $ & $R_{2,k-2}^{ }$ & $R_{2,k-1}^{ }$ & \\
			3      &  $R_{3,0}^{ }$ & $R_{3,1}^{ }$ & $R_{3,2}^{ }$ & $ \dots $ & $R_{3,j}^{ }$ & $ \dots $ & $R_{3,k-2}^{ }$ & & \\
			4      &  $R_{4,0}^{ }$ & $R_{4,1}^{ }$ & $R_{4,2}^{ }$ & $ \dots $ & $R_{4,j}^{ }$ & $ \dots $ & & & \\
			:      & & & & & & & & & \\
			i      &  $R_{i,0}^{ }$ & $R_{i,1}^{ }$ & $R_{i,2}^{ }$ & $ \dots $ & $R_{i,j}^{ }$ & & & & \\
			:      & & & & & & & & & \\
			s-2      &  $R_{s-2,0}^{ }$ & $R_{s-2,1}^{ }$ & $R_{s-2,2}^{ }$ & & & & & & \\
			s-1      &  $R_{s-1,0}^{ }$ & $R_{s-1,1}^{ }$ & & & & & & & \\
			s      &  $R_{s,0}^{ }$ & & & & & & & & \\
		\end{tabularx}
	\caption{Matriz de Residuales Brutos}
	\label{matrix5}
	\end{table}
 
\section{Aplicación del Método de Bootstrap}
 
	 \citet{QLowe} indica que es correcto realizar un muestreo con reemplazo de residuales n veces sobre el triangulo de siniestralidad ajustada mediante  la seleccion aleatoria de entre los residuales obtendidos. Sin embargo, con el fin de simular un comportamiento realista acorde a la información historica usada, se emplea el método de $Bootstrap$, bajo el supuesto de que los residuales brutos $R_{i,j}$ de la matriz ~\ref{matrix5} provienen de la misma distribución y son independientes, y se definen $R_{j}^{\text{min}}$ como el valor mínimo de los residuales observados en el trimestre reportado $j$ y $R_{j}^{\text{max}}$ como el valor máximo de los residuales observados en el trimestre reportado $j$, como se observa en la tabla ~\ref{matrix6}.
	
	\begin{table}[ht]
		\centering
		\begin{tabularx}{\linewidth}{ c|ccccccccc}
			& \multicolumn{9}{c}{Trimestre en que se reportó el procedimiento} \\
			& 0  & 1 & 2 & $ \dots $ & j & $\dots $ & k-2 & k-1 &  k\\
			\midrule
			Mínimo      &  $R_{0}^{\text{min}}$ & $R_{1}^{\text{min}}$ & $R_{2}^{\text{min}}$ & $ \dots $ & $R_{j}^{\text{min}}$ & $ \dots $ & $R_{k-2}^{\text{min}}$ & $R_{k-1}^{\text{min}}$ & $R_{k}^{\text{min}}$ \\
			Máximo      &  $R_{0}^{\text{max}}$ & $R_{1}^{\text{max}}$ & $R_{2}^{\text{max}}$ & $ \dots $ & $R_{j}^{\text{max}}$ & $ \dots $ & $R_{k-2}^{\text{max}}$ & $R_{k-1}^{\text{max}}$ & $R_{k}^{\text{max}}$ \\
		\end{tabularx}
	\caption{Residuales Mínimos y Máximos de cada trimestre}
	\label{matrix6}
	\end{table}
	
Donde
	
\begin{itemize}
	
	\setlength\itemsep{-0.5em}
	
	\item $R_{j}^{\text{min}}= \text{min}_{ i\in \left\{1,2,\dots ,S\right\}} \left[R_{i,j}^{}\right]$	
	
	\item $R_{j}^{\text{max}}= \text{max}_{ i\in \left\{1,2,\dots ,S\right\}} \left[R_{i,j}^{}\right]$
	
	\item $R_{i,j}$ = Residual bruto del trimestre de inicio de vigencia $i$ reportado en el trimestre $j$.
		
\end{itemize} 

	De este modo se seleccionan los valores de los residuales de entre estos mínimos y máximos, realizando un muestreo con reemplazo de residuales de $n$ muestras de forma uniforme dentro del intervalo $\left[R_{j}^{\text{min}},R_{j}^{\text{max}}\right]$ de cada uno de los trimestres de reporte $j$. Se obtiene la matriz de residuales (tabla ~\ref{matrix7}) definiendo $R_{i,j}^{*}$ como el residual de la muestra de pólizas con inicio de vigencia en el trimestre $i$ con trimestre de reporte $j$. 
		
	\begin{table}[H]
		\centering
		\begin{tabularx}{\linewidth}{ c |ccccccccc}
			\multirow{2}{4cm}{Trimestre de inicio de vigencia de la póliza}& \multicolumn{9}{c}{Trimestre en que se reportó el procedimiento} \\
			& 0  & 1 & 2 & $ \dots $ & j & $\dots $ & k-2 & k-1 &  k\\
			\midrule
			1      &  $R_{1,0}^{*}$ & $R_{1,1}^{*}$ & $R_{1,2}^{*}$ & $ \dots $ & $R_{1,j}^{*}$ & $ \dots $ & $R_{1,k-2}^{*}$ & $R_{1,k-1}^{*}$ & $R_{1,k}^{*}$ \\
			2      &  $R_{2,0}^{*}$ & $R_{2,1}^{*}$ & $R_{2,2}^{*}$ & $ \dots $ & $R_{2,j}^{*}$ & $ \dots $ & $R_{2,k-2}^{*}$ & $R_{2,k-1}^{*}$ & \\
			3      &  $R_{3,0}^{*}$ & $R_{3,1}^{*}$ & $R_{3,2}^{*}$ & $ \dots $ & $R_{3,j}^{*}$ & $ \dots $ & $R_{3,k-2}^{*}$ & & \\
			4      &  $R_{4,0}^{*}$ & $R_{4,1}^{*}$ & $R_{4,2}^{*}$ & $ \dots $ & $R_{4,j}^{*}$ & $ \dots $ & & & \\
			:      & & & & & & & & & \\
			i      &  $R_{i,0}^{*}$ & $R_{i,1}^{*}$ & $R_{i,2}^{*}$ & $ \dots $ & $R_{i,j}^{*}$ & & & & \\
			:      & & & & & & & & & \\
			s-2      &  $R_{s-2,0}^{*}$ & $R_{s-2,1}^{*}$ & $R_{s-2,2}^{*}$ & & & & & & \\
			s-1      &  $R_{s-1,0}^{*}$ & $R_{s-1,1}^{*}$ & & & & & & & \\
			s      &  $R_{s,0}^{*}$ & & & & & & & & \\
		\end{tabularx}
		\caption{Matriz de Residuales}
		\label{matrix7}
	\end{table}

	\vspace{0.1cm}
	
	Se obtiene así una matriz de siniestros simulada al agregar los residuales a cada monto de siniestros ajustados (tabla ~\ref{matrix8}), definiendo $X_{i,j}^{sim}$ como el monto de siniestros simulado de las pólizas con inicio de vigencia en el trimestre $i$ que fueron reportados $j$ trimestres posteriores al inicio de vigencia (ver eq. \ref{eq8}).
	
	\begin{equation}
    X_{i,j}^{sim}=R_{i,j}^{*}+X_{i,j}^{*} 
\label{eq8}    
\end{equation}
	
	\vspace{0.1cm}
	
	Donde
	
	\vspace{0.1cm}
	
	\begin{itemize}
		\setlength\itemsep{-0.5em}
		\item $R_{i,j}^{*}$ = Residual seleccionado en la muestra que corresponde a las pólizas con inicio de vigencia en el trimestre $i$ con trimestre de reporte $j$.
		\item $X_{i,j}^{*}$ = Monto de siniestros ajustados de las pólizas con inicio de vigencia en el trimestre $i$ que fueron reportados $j$ trimestres posteriores al inicio de vigencia.
		\end{itemize}
	
	
	\begin{table}[H]
		\centering
		\begin{tabularx}{\linewidth}{ c|ccccccccc}
			\multirow{2}{4cm}{Trimestre de inicio de vigencia de la póliza}& \multicolumn{9}{c}{Trimestre en que se reportó el procedimiento} \\
			& 0  & 1 & 2 & $ \dots $ & $j$ & $\dots $ & k-2 & k-1 &  k \\
			\midrule
			1      &  $X_{1,0}^{sim}$ & $X_{1,1}^{sim}$ & $X_{1,2}^{sim}$ & $ \dots $ & $X_{1,j}^{sim}$ & $ \dots $ & $X_{1,k-2}^{sim}$ & $X_{1,k-1}^{sim}$ & $X_{1,k}^{sim}$ \\
			2      &  $X_{2,0}^{sim}$ & $X_{2,1}^{sim}$ & $X_{2,2}^{sim}$ & $ \dots $ & $X_{2,j}^{sim}$ & $ \dots $ & $X_{2,k-2}^{sim}$ & $X_{2,k-1}^{sim}$ & \\
			3      &  $X_{3,0}^{sim}$ & $X_{3,1}^{sim}$ & $X_{3,2}^{sim}$ & $ \dots $ & $X_{3,j}^{sim}$ & $ \dots $ & $X_{3,k-2}^{sim}$ & & \\
			4      &  $X_{4,0}^{sim}$ & $X_{4,1}^{sim}$ & $X_{4,2}^{sim}$ & $ \dots $ & $X_{4,j}^{sim}$ & $ \dots $ & & & \\
			:      & & & & & & & & &\\
			i      &  $X_{i,0}^{sim}$ & $X_{i,1}^{sim}$ & $X_{i,2}^{sim}$ & $ \dots $ & $X_{i,j}^{sim}$ & & & & \\
			:      & & & & & & & & & \\
			s-2      &  $X_{s-2,0}^{sim}$ & $X_{s-2,1}^{sim}$ & $X_{s-2,2}^{sim}$ & & & & & & \\
			s-1      &  $X_{s-1,0}^{sim}$ & $X_{s-1,1}^{sim}$ & & & & & & & \\
			s      &  $X_{s,0}^{sim}$ & & & & & & & & \\
		\end{tabularx}
	\caption{Matriz de Siniestros Simulados}
	\label{matrix8}
	\end{table}

Usando esta matriz de siniestros simulada, se genera la matriz de siniestros acumulados simulados (como se generó la tabla ~\ref{matrix2}), se obtiene los factores de incremento simulados (como se obtuvieron en la eq. \ref{eq2}), se estiman los siniestros esperados simulados (como se estimaron en la eq. \ref{eq3}) y se calculan los flujos de obligaciones futuras simuladas de riesgos en curso de la muestra $i$ ($RRC_{i}^{sim}$) (como se calcularon en la eq. \ref{eq4}), repitiendo el proceso $n$ veces.
	
	 
	
	Se puede considerar el mejor estimador de riesgos en curso ($BELR_{RRC}$), como el valor medio de las $n$ muestras de los flujos de obligaciones futuras simuladas de riesgos en curso (ver eq \ref{eq9}).


	 

	
	\begin{equation}
		BELR_{RRC}^{}=\frac{\sum _{i=1}^{n}RRC_{i}^{sim}}{n}
		\label{eq9}
		\end{equation}
			


 
	
	Donde
	
	 
		\begin{itemize}
		\setlength\itemsep{-0.5em}
	\item $RRC_{i}^{sim}=$ i-ésima simulación de los flujos de obligaciones futuras de riesgos en curso.
	
	\item $n=$ número de simulaciones realizadas.
	\end{itemize}
	
\begin{comment}	
	 
	Para Salud Individual Dental, el $BELR_{RRC}$ se calculará como el producto de la prima de tarifa no devengada y el factor de siniestralidad última de mercado proporcionado por la Comisión Nacional de Seguros y Fianzas:


	 

$ $

 
	
	
	{\centering
		$BELR_{RRC}^{}=PTND_{} \cdot FS_{BEL}^{RRC}$
		\noindent
		
	}
	
	 

$ $

 
	
	Donde
	
	 
	
	$PTND_{}=$ Prima de tarifa no devengada.
\begin{comment}	
	$FS_{BEL}^{RRC}=$ Factor de Siniestralidad última con información de mercado.
	
	 

$ $

 

	El cálculo del $BELR_{RRC}$, se realizará de forma trimestral y se prorrateara con el vigor de la valuación del cierre de mes a fin de obtener la reserva de riesgos en curso.

\end{comment}
	
	 
\chapter{Reserva de Gastos de Administración}\label{metnum}

Para el cálculo de los componentes faltantes de la Reserva de Riesgos en Curso, así como para el prorrateo del $BELR_{RRC}$ por asegurado, es necesario primero obtener un Factor de No Devengamiento, con el cual se podrá calcular la Reserva de Gastos y la Prima de Riesgos No Devengada.

\section{Cálculo del Factor de No Devengamiento}

	El Factor de No Devengamiento ($F_{ND}$) de cada asegurado en vigor equivale a la proporción de tiempo no transcurrido de la vigencia del asegurado (ver eq. \ref{eq13}).


\begin{equation}	
	F_{ND}=\begin{cases}
		0, & \text{si } FVal \geqslant FFin.\\
		\frac{FFin-FVal}{FFin-FIni}, & \text{si } FIni \leqslant FVal \leqslant FFin\\
		1, & \text{si } FIni \geqslant FVal
	\end{cases}
	\label{eq13}
\end{equation}


$ $

Donde


\begin{itemize}
	\setlength\itemsep{-0.5em}
	\item $FIni=$ Fecha de Inicio de cobertura para el asegurado.
	
	\item $FFin=$ Fecha de Fin de cobertura para el asegurado.
	
	\item $FVal=$ Fecha de Valuación.
	
\end{itemize}



\section{Cálculo del Bel de Gastos de Administración}

Se debe calcular ahora la reserva de gastos de administración no devengados. Esta es la parte no devengada correspondiente a la porción de prima de tarifa anual de cada una de las pólizas en vigor al momento de la valuación\footnote{Ver \citet{SAnexoCUSF}, Apartado A2.1.7.4}.

Para realizar el cálculo es necesario identificar primero el número de asegurados en vigor al cierre del mes al momento de la valuación, el monto de prima correspondiente a los beneficios contratados, los gastos asociados y el periodo de cobertura de cada asegurado en vigor.



El  Monto de los Gastos de Administración no devengado de cada asegurado ($BELG_{ADM,ind}$) se define entonces como la porción correspondiente a los Gastos de Administración de la Prima de Tarifa de los asegurados de que se trate multiplicada por el Factor de No Devengamiento (ver eq. \ref{eq10}). 


\begin{comment}
Para Salud Individual Dental, se tomara el porcentaje de gasto de administración del mercado proporcionado por la Comisión Nacional de Seguros y Fianzas.
\end{comment}	



\begin{comment}	
$
BELG_{ADM,ind}=\begin{cases}
PT \cdot G_{ADM} \cdot F_{ND}, & \text{$g \neq Salud  Dental  Individual$}.\\
PT \cdot \alpha_{i} \cdot F_{ND}, & \text{$g = Salud  Dental  Individual$}.
\end{cases}
$
\end{comment}

\begin{equation}
	BELG_{ADM,ind}=	PT \cdot G_{ADM} \cdot F_{ND}
	\label{eq10}
\end{equation}



Donde

\begin{itemize}
	\setlength\itemsep{-0.5em}
	\item $PT=$ Prima de tarifa
	
	\item $G_{ADM}^{}=$ Gasto de Administración.
	
	\item $F_{ND}^{}=$ Factor de No Devengamiento.
\end{itemize}

\begin{comment}	
$\alpha_{i}^{}=$ Porcentaje de gasto de administración con información de Mercado
\end{comment}	


Se define el mejor estimador de la reserva de Gastos de Administración ($BELG_{ADM}$) como el monto integrado por la suma de los Gastos de Administración no devengados de cada uno de los asegurados en vigor (ver eq. \ref{eq11}).


\begin{equation}
	BELG_{ADM}=\sum _{}^{}BEL_{ADM,ind}^{}
	\label{eq11}
\end{equation}



	\section{Cálculo de la Prima de Riesgo No Devengada}


La Prima de Riesgo No Devengada corresponderá al valor de la prima de riesgo multiplicada por el factor de No Devengamiento.

Para realizar el cálculo se determinará para cada uno de los asegurados en vigor su Prima de Riesgo correspondiente. Ya que la Prima de Tarifa equivale a la Prima de Riesgo más gastos de adquisición, más gastos de administración, más un gasto de utilidad\footnote{Ver \citet{TZurich}, Prima de Tarifa}., es posible deducir el valor de la Prima de Riesgo con un cálculo sencillo (ver eq. \ref{eq12}).


\begin{equation}
	{PR}_{}^{}={{PT}_{}\cdot(1-G_{ADM}-C_{ADQ}-U)}
	\label{eq12}
\end{equation}


Donde


\begin{itemize}
	\setlength\itemsep{-0.5em}
	\item $PT=$ Prima de Tarifa.
	
	\item $G_{ADM}^{}=$ Gasto de Administración.
	
	\item $C_{ADQ}^{}=$ Costo de Adquisición.
	
	\item $U_{}^{}=$ Margen de Utilidad.
\end{itemize}


Una vez que se tiene la Prima de Riesgo se puede calcular la Prima de Riesgo no Devengada de cada uno de los asegurados ($PRND_{ind}$), dada como la Prima de Riesgo multiplicada por el Factor de No Devengamiento calculado en la eq. \ref{eq13} (ver eq. \ref{eq14}).	


\begin{equation}
	{PRND}_{ind}^{}={{PR}_{ind}\cdot F_{ND}}
	\label{eq14}
\end{equation}



Donde

\begin{itemize}
	\setlength\itemsep{-0.5em}	 
	
	\item $F_{ND}=$ Factor de No Devengamiento de cada asegurado.
	
	\item $PR_{ind}=$ Prima de Riesgo de cada asegurado.
\end{itemize}


\begin{comment}

Mientras que para las pólizas emitidas anticipadamente que, al momento de la valuación, no han iniciado vigencia, la prima de riesgo no devengada se calculará como la diferencia entre la Prima de Tarifa de cada asegurado y el Monto de Gastos de Administración no devengado de cada asegurado (ver eq. \ref{eq15}).



\begin{equation}
	{PRND}_{ind}^{}={{PT}_{ind}-BELG_{ADM,ind}}
	\label{eq15}
\end{equation}


Donde


\begin{itemize}
	\setlength\itemsep{-0.5em}
	\item $PT_{ind}=$ Prima de Tarifa de cada asegurado.
	
	\item $BELG_{ADM,ind}=$ Monto de los gastos de administración no devengado de cada asegurado.
\end{itemize}
\end{comment}

Así pues, la prima de riesgo no devengada ($PRND_{}$) se calculará como la suma de los montos de Prima de riesgo no devengada de cada asegurado (ver eq. \ref{eq16}).

\begin{equation}
	PRND_{}={\sum _{}^{}PRND_{ind}^{}}
	\label{eq16}
\end{equation}


Donde $PRND_{ind}=$ Prima de riesgo no devengada de cada asegurado.


	 
	 	\chapter{Cálculo de la Reserva de Riesgos en Curso}\label{metnum}
	 
	 Ya que previamente se calculó el Mejor Estimador de Riesgos en curso ($BELR_{RRC}$) del grupo homogéneo, este debe ser prorrateado para cada asegurado en vigor. Para esto es necesario obtener un Factor de Distribución ($FD_{}$) que usa la $PRND_{}$ para realizar el prorrateo de acuerdo a la proporsión de Prima de Riesgo No Devengada que corresponde a cada asegurado ($PRND_{ind}$) (ver eq. \ref{eq17}).
	 
	 
	 \begin{equation}
	 	{FD}_{}^{}=\frac{{BELR}_{RRC}^{}}{{PRND}_{}}
	 	\label{eq17}
	 \end{equation}
	 
	 
	 Donde
	 
	 
	 \begin{itemize}
	 	\setlength\itemsep{-0.5em}
	 	\item $BELR_{RRC}=$ Mejor Estimador de Riesgos en Curso.
	 	
	 	\item $PRND_{}=$ Prima de Riesgo No Devengada calculada previamente.
	 	
	 \end{itemize}
	 
	 $ $
	 
	 
	 Así, el mejor estimador de riesgos en curso individual ($BELR_{RRC,ind}$) se obtiene multiplicando la Prima de Riesgo No Devengada de cada asegurado por el Factor de Distribución (ver eq. \ref{eq18}).
	 
	 
	 
	 \begin{comment}	
	 $
	 BELR_{RRC,ind}=\begin{cases}
	 PRND_{ind}\cdot FD_{}, & \text{$g \neq Salud Dental Individual$}\\
	 
	 PTND_{ind}\cdot FS_{BEL}^{RRC}, & \text{$g = Salud Dental Individual$}
	 \end{cases}
	 $
	 \end{comment}	
	 
	 \begin{equation}
	 	BELR_{RRC,ind}=PRND_{ind}\cdot FD_{}
	 	\label{eq18}
	 \end{equation}
	 
	 
	 
	 Donde
	 
	 
	 \begin{itemize}
	 	\setlength\itemsep{-0.5em}	
	 	\item $PRND_{ind}=$ Prima de Riesgo no Devengada de cada asegurado.
	 	
	 	\item $FD_{}^{}=$ Factor de Distribución.
	 \end{itemize}
	 
	 \begin{comment}
	 $PTND_{ind}=$ Prima de Tarifa no Devengada de cada asegurado
	 
	 
	 $FS_{BEL}^{RRC}=$ Factor de Siniestralidad última con información de mercado
	 \end{comment}
	 
	 

 
	

	
	

 
	
	\section{Cálculo de la Reserva de Riesgos en Curso de cada Asegurado}
	
	 
	Ya que se tienen todos los componentes, la Reserva de Riesgos en Curso de cada asegurado ($RRC_{ind}$) resulta de sumar el Mejor Estimador de Riesgos en Curso de cada asegurado con el Monto de Gastos de Administración de cada asegurado (ver eq. \ref{eq19}).
	
\begin{equation}
		{RRC}_{ind}^{}={{BELR}_{RRC,ind}+BELG_{ADM,ind} +MR_{RRC,ind}
		}
	\label{eq19}
		\end{equation}	
			
	Donde
	 
		\begin{itemize}
		\setlength\itemsep{-0.5em}
	\item $BELR_{RRC,ind}=$ Mejor estimador de riesgos en curso individual.
	
	\item $BELG_{ADM,ind}^{}=$ Monto no devengado de Gastos de Administración de cada asegurado.
	
		\item $MR_{RRC,ind}^{}=$ Margen de Riesgo de la reserva de riesgos en curso calculado usando el RCS del cierre del mes inmediato anterior a la fecha de valuación.
	\end{itemize}
	



		
	Por lo tanto, la Reserva de Riesgos en Curso Total ($RRC_{}$) se calculará como la suma del Monto de Reserva de Riesgos en Curso de cada asegurado (ver eq. \ref{eq20}).
	
\begin{equation}
		RRC_{}={\sum _{}^{}RRC_{ind}^{}}
		\label{eq20}
	\end{equation}
		
	
	Donde $RRC_{ind}^{}=$ Reserva de Riesgos en Curso de cada asegurado.
	


 

%------------------------------------------------------------------------------------------ %
\clearpage
\appendix 

\fancyhead[LE]{\scshape\thepage\hspace{1cm}\footnotesize\nouppercase{Apéndices}}
\fancyhead[RO]{\scshape\footnotesize\nouppercase{Apéndice \thechapter}\hspace{1cm}\normalsize\thepage}
\fancyhead[LO]{}
\fancyhead[RE]{}
\pagestyle{fancy}

\chapter{Metodología codificada en R }\label{desarrolloconstadm}

A continuación se presenta la Metodología descrita en este trabajo convertida a codigo de R. Se realiza con el fin de corroborar los resultados de la Metodología aplicada manualmente, así como facilitar las $n$ simulaciones necesarias.


\begin{lstlisting}
## Se instalan los paquetes necesarios de R si aún no se tienen instalados
install.packages("ChainLadder")
install.packages("psych")

## Se referencían los paquetes instalados mediante las bibliotecas
library(ChainLadder)
library(psych)

## Función Remover para limpiar la consola
rm(list=ls())

## Se define la función con la cual se hara el Cálculo de Reservas RRC
SIMULACION_RRC <- function (MATRIX_INICIAL, n, archivo){
	
	## Se define la dimensión (m) de los triángulos de residuales iniciales
	library(ChainLadder)
	m=dim(MATRIX_INICIAL)[1]
	
	
	## Inicia el cálculo
	
	## La siguiente sección corresponde a los datos obtenidos mediante la plantilla de Excel
	
	## Mediante la función "incr2cum", la cual transforma un triángulo incremental en triángulo acumulativo, donde MATRIX_INICIAL se refiere al triángulo inicial de siniestros como muestra la tabla 5.1, se genera MATRIX_SIN_ACUM, que es el triángulo de sinietros acumulados como muestra la tabla 5.2
	MATRIX_SIN_ACUM <- incr2cum(MATRIX_INICIAL)
	
	## Se obtienen los factores fj como indica la eq. 5.2, usando la matriz de siniestros acumulados
	FACTORES_J <- sapply(1:(m-1), function(i){ sum(MATRIX_SIN_ACUM[c(1:(m-i)), i+1])/sum(MATRIX_SIN_ACUM[c(1:(m-i)), i]) } )
	
	## Mediante estos factores de incremento, se definen los Siniestros Esperados como indica la eq. 5.3
	SIN_ESPERADOS <- cbind(MATRIX_SIN_ACUM, RVA = rep(0,m))
	for (j in 1:m){
		for(i in 1:m){
			if(is.na(SIN_ESPERADOS[i,j])){
				SIN_ESPERADOS[i,j] = SIN_ESPERADOS[i,(j-1)]*FACTORES_J[j-1]
			}
		}
	}
	
	## Con los Siniestros Esperados, se definen las obligaciones futuras iniciales de riesgos en curso RRC0 como indica la eq. 5.4 
	for (i in 1:m){
		SIN_ESPERADOS[i,(m+1)] = SIN_ESPERADOS[i,m]-SIN_ESPERADOS[i,(m-i+1)]
	}
	
	## Se obtiene el monto ajustado de siniestros de las pólizas como indica la eq. 5.5 y con ellos se genera el triángulo de siniestros acumulados ajustados como se observa en la tabla 5.3
	MATRIX_ACUM_AJUST = MATRIX_SIN_ACUM
	for (i in 1:(m-1)){
		for(j in 1:(m-i)){
			MATRIX_ACUM_AJUST[i,(m-j-i+1)] = MATRIX_ACUM_AJUST[i,(m-j-i+2)]/FACTORES_J[m-i-j+1]
		}
	}
	
	## A partir de este triángulo se obtiene el triengulo de siniestros ajustados como muestra la tabla 5.4 usando la función "cum2incr", la cual transforma un triángulo acumulativo en un triángulo incremental
	MATRIX_SIN_AJ = cum2incr(MATRIX_ACUM_AJUST)
	
	## Se obtienen los Residuales Brutos como indica la eq. 5.7
	RESIDUAL = MATRIX_INICIAL - MATRIX_SIN_AJ
	
	## Se define una nueva Matriz de n x m-1 para las n simulaciones
	SIMULACION_RESERVA=c()
	
	## Inician las n simulaciones
	for (k in 1:n){
		# Se define la Muestra de Residuales con dimensiones n x m
		MUESTRA_RESIDUALES = matrix(nrow=m, ncol=m)
		
		## Se definen el minimo y maximo de cada trimestre como se observa en la tabla 5.6 y se genera el triángulo de residuales como se observa en la tabla 5.7
		for(j in 1:m){
			MUESTRA_RESIDUALES[,j] = t(runif(dim(RESIDUAL)[1],min=min(na.omit(RESIDUAL[,j])), max=max(na.omit(RESIDUAL[,j]))))
		}
		
		## Se agregan el residual obtenido a cada monto de siniestros ajustados como indica la eq. 5.8 para obtener los siniestros simulados y con esto se obtiene el triángulo de Siniestros simulados como se ve en la tabla 5.8		
		SINIESTRO_SIM = MUESTRA_RESIDUALES + MATRIX_SIN_AJ
		
		## Inicia la repetición del proceso original
		
		## Aplicando la función "incr2cum" al triángulo obtenido se transforma el triángulo incremental a un triángulo de siniestros acumulados simulados como se genero originalmente la tabla 5.2
		MATRIX_SIN_SIM <- incr2cum(SINIESTRO_SIM)
		
		## Se obtinenen nuevamente los factores de incremento como indica la eq. 5.2
		FACTOR_N_SIM <- sapply(1:(m-1), function(i){ sum(MATRIX_SIN_SIM[c(1:(m-i)), i+1])/sum(MATRIX_SIN_SIM[c(1:(m-i)), i]) } )
		
		## Se estiman los siniestros esperados simulados como indica la eq. 5.3
		SIN_ESP_N_SIM <- cbind(MATRIX_SIN_SIM, RVA = rep(0,m))
		for (j in 1:m){
			for(i in 1:m){
				if(is.na(SIN_ESP_N_SIM[i,j])){
					SIN_ESP_N_SIM[i,j] = SIN_ESP_N_SIM[i,(j-1)]*FACTOR_N_SIM[j-1]
				}
			}
		}
		
		## Se calculan los flujos de obligaciones futuras simuladas de riesgos en curso como indica la eq. 5.4
		for (i in 1:m){
			SIN_ESP_N_SIM[i,(m+1)] = SIN_ESP_N_SIM[i,m]-SIN_ESP_N_SIM[i,(m-i+1)]
		}
		
		## La reserva calculada corresponde a la suma total de los resultados
		SIMULACION_RESERVA[k]=sum(SIN_ESP_N_SIM[,m+1])
		
		
	}
	
	## Se definen las Obligaciones Futuras Iniciales de Riesgos en Curso inicial como la primera iteracion del metodo
	RRC_INICIAL=sum(SIN_ESPERADOS[,m+1])
	
	## Se define el Mejor Estimador de Riesgos en Curso como el valor medio de las n muestras de los flujos de obligaciones futuras sumuladas como indica la eq. 5.9
	RRC_MEDIA=mean(SIMULACION_RESERVA)
	
	
	## Ya que se terminan las n simulaciones, para obtener los resultados en Excel se emplea la funcion "write.table" para que arroje las cifras de las n simuladciones en un archivo de texto
	write.table (SIMULACION_RESERVA, file = paste("RRC_N_SIM_", archivo, ".txt", sep=""), quote = FALSE, sep = "\t", eol = "\n", dec = ".", row.names = FALSE, col.names = FALSE)
	print(paste("RRC_N_SIM_", archivo, ".txt", sep=""))
	
	## Ya que se terminan las n simulaciones, para obtener los resultados del proceso en R se emplea la funcion "write.table" para que arroje el valor de la primera iteracion y la media de las n simulaciones
	write.table (c(RRC_INICIAL=RRC_INICIAL, media=RRC_MEDIA), file = paste("RRC_INICIAL_MEDIA", archivo, ".txt", sep=""), quote = FALSE, sep = "\t", eol = "\n", dec = ".", row.names = FALSE, col.names = FALSE)
	print(paste("RRC_INICIAL_MEDIA", archivo, ".txt", sep=""))
	
	
	print("FIN DE LA FUNCION")
	
	## Repetir el proceso hasta que se terminen las n simulaciones
	return(list(SIMULACION_RESERVA, c(RRC_INICIAL=RRC_INICIAL, media=RRC_MEDIA)))
	
}


## Aplicar la metodologia en R al triángulo de siniestros guardado en formato .csv en la ruta correspondiente
triangulo_SINIESTROS_CSV<-read.csv("D:/RESPALDO/Tesis Chris/triangulo Siniestros Prueba.csv")
RESULTADO_RRC<-SIMULACION_RRC(triangulo_SINIESTROS_CSV, 100000, "triangulo_SINIESTROS_CSV")


\end{lstlisting}

	
	\pagestyle{empty}
\bibliography{library}
	
\end{document}
